\section{Conditional counts and discrete curvature}\label{ms:counts}
The sign decomposition turns an iid sum into a binomial mixture of
conditional noise sums. This identity supplies a comparison with the
Bernoulli law having the same sign probability. The main issue is to
control the error proportionally to the noise variance, including when
that variance is arbitrarily small. Throughout this section the target is
$c=\CE$, and $a$ is the fixed-target excess. The identities below hold for
any actual marginal with the stated moments; extremality is used only when
we apply them to a witness with $a\ge0$.

\subsection{The exact count and score identities}

Let $\pi_j=\binom njp^jq^{n-j}$ for $0\le j\le n$, zero otherwise,
and $F_B$ be the CDF of $\operatorname{Bin}(n,p)$. For independent copies
of the marginal, conditioning on the entire sign vector gives independent
conditional residuals. All sign vectors with $j$ positive entries yield
the same sum law
\[
 V_j=U_{1,1}+\cdots+U_{1,j}+U_{0,j+1}+\cdots+U_{0,n}.
\]
It follows from \eqref{ms:bernoulli-coordinates} that at
$x=np+H/\ell=k+u$,
\begin{equation}
 \Pp(S_n\le H)=\sum_{j=0}^n\pi_j\Pp(V_j\le k+u-j).
\label{ms:count-identity}
\end{equation}
For $0\le k\le n$, subtracting $F_B(k)$ gives the exact signed identity
\begin{align}
 Q:={}&\Pp(S_n\le H)-F_B(k)\notag\\
   ={}&-\sum_{i=0}^{k}\pi_{k-i}\Pp(V_{k-i}>u+i)
      +\sum_{i=1}^{n-k}\pi_{k+i}\Pp(V_{k+i}\le u-i).
\label{ms:count-difference}
\end{align}
The outgoing event is strict and the incoming events are closed. Thus an
atom exactly on a count plane has its correct weight.

Choosing a nearest integer gives $|u|\le1/2$, with either integer allowed
at a tie. Restricting to legal counts $0\le k\le n$ requires either
$-1/2\le np+H/\ell\le n+1/2$ or a separate treatment of the remaining
thresholds. Interior-count normalizations further require $1\le k\le n-1$.

The own-count variance is
\[
 t=\E V_k^2=(n-k)w_0+kw_1=\lambda+\delta(w_1-w_0),\qquad
 \delta=k-np=H/\ell-u.
\]
Since $|w_1-w_0|\le s/\min(p,q)$, any established bound
$|\delta|\le\Delta$ and $c_v\ge1/\min(p,q)$ gives, for $n\ge N$,
\begin{equation}
 (1-c_v\Delta/N)\lambda\le t\le(1+c_v\Delta/N)\lambda.
\label{ms:own-count-variance}
\end{equation}
When the lower bound is used in a denominator, $1-c_v\Delta/N$ must be positive.

For the comparison with the same Bernoulli count, set
\[
 \rho_B=(1-2v)/\sqrt v,\qquad
 C=\frac{\sqrt n}{\rho}\{\Pp(S_n\le H)-\Phi(H/\sqrt n)\},
\]
\[
 C_B=\frac{\sqrt n}{\rho_B}
       \{F_B(k)-\Phi(\delta/\sqrt{nv})\},
 \qquad
 f(u)=\frac{\Phi((\delta+u)/\sqrt{n(v+s)})-
                    \Phi(\delta/\sqrt{n(v+s)})}{\pi_k}.
\]
For a legal count $\pi_k>0$. Define the exact repair
\[
 \mathcal R=\sqrt n\{\Phi(\delta/\sqrt{nv})-
                    \Phi(\delta/\sqrt{n(v+s)})\}
                         +c(\rho_B-\rho).
\]
Direct subtraction, using $H/\sqrt n=(\delta+u)/\sqrt{n(v+s)}$, yields
\begin{equation}
 a+\rho_B(c-C_B)=\sqrt n\pi_k\{Q/\pi_k-f(u)\}+\mathcal R.
\label{ms:master-identity}
\end{equation}
Both terms on the right must be estimated for the same law, threshold and phase.

\begin{proposition}[Two conditional local witnesses]
\label{ms:local-witnesses}
Let $\lambda>0$, $n\ge N$, and suppose $\sqrt n\pi_k\ge m_0>0$.
The following are sufficient for $a+\rho_B(c-C_B)<-\varepsilon\lambda$.
\begin{enumerate}
\item Bounds $O,T,R_0,\kappa$ satisfy
\[
 \Pp(V_k>u)+f(u)\ge O,\quad
 \sum_{i\ge1}\Pp(V_{k+i}\le u-i)\le T,\quad
 \pi_{k+i}/\pi_k\le R_0,\quad
 \mathcal R\le\kappa\lambda/n,
\]
where the sum and ratio involve legal counts. If $D_0=O-R_0T>0$ and
$m_0D_0-\kappa\lambda/n>\varepsilon\lambda$, the conclusion holds.
\item A proved whole-count estimate gives
\[
 a+\rho_B(c-C_B)\le\sqrt n\pi_k\lambda B(\lambda)
                               +\lambda I(\lambda)/n.
\]
If $B(\lambda)<0$ and
$m_0B(\lambda)+\max(0,I(\lambda))/N<-\varepsilon$, the conclusion holds.
\end{enumerate}
\end{proposition}
\begin{proof}
For the first witness, discard only the outgoing terms with $i>0$ in
\eqref{ms:count-difference}; then $Q/\pi_k-f(u)\le-D_0$.
Use $D_0>0$ before replacing the count mass by its lower bound in
\eqref{ms:master-identity}. For the second, $B<0$ similarly licenses the
mass replacement, and $I/n\le\max(0,I)/N$ holds for either sign of $I$.
Multiplication by $\lambda>0$ gives the asserted strict inequality.
\end{proof}

For a witness with $a\ge0$, Schulz's theorem $C_B\le c$ makes
$a+\rho_B(c-C_B)$ nonnegative. Either strict negative bound therefore
excludes that witness~\cite{Schulz2016}. This is a comparison of actual
iid sums; it is not an extremal principle for a moment relaxation.

At zero noise the threshold may be taken at a lattice point: at an attained
positive maximum for a Bernoulli law, the threshold is a binomial atom. Between consecutive atoms, lowering the threshold
to the preceding atom preserves the CDF and strictly improves the normal
subtraction. Above the final atom the same argument applies, while below
the first atom the discrepancy is nonpositive. Thus $u=0$ and $C=C_B$
at such a maximum. The sharp Bernoulli upper bound gives $a\le0$.
To exclude equality as well, one needs its strict finite-index form; we
keep that requirement explicit in the conditional global statement.

\subsection{Whole-count coupling and its curvature term}

Part of the whole-count estimate follows from a coupling. Suppose
$U_0\le a_0$, $U_1\ge-a_1$, where $d_Y=a_0+a_1<1/2$, and put $\zeta=1-d_Y$.
Independently realize all $2n$ variables in the definition of $V_j$ above,
and extend $V_j$ constantly for $j<0$ and $j>n$. The increments of
$j+V_j$ are at least $\zeta>1/2$ inside the legal range and equal one outside.
Consequently
\[
 K=\max\{j\in\mathbb Z:j+V_j\le k+u\},\qquad R=K-k
\]
are well defined. An independent binomial selector has the correct law
for each count, so
\begin{equation}
 \Pp\Bigl(\sum_{i=1}^nY_i\le k+u\Bigr)=\E F_B(k+R).
\label{ms:whole-count-cutoff}
\end{equation}
Indeed, conditional on the coupled residuals, the closed event for the
selector occurs exactly at counts not exceeding $K$.

\begin{lemma}[Discrete curvature and crossing moments]
\label{ms:crossing-curvature}
Let $\chi=(n+1)p-k<2$ and $m_B=\max_j\pi_j$. For every integer $r$,
\begin{align}
 F_B(k+r)-F_B(k)-r\pi_k
 \le\frac{m_B}{(n+1)v}
 \left\{(\chi-1)_+r_++\chi_+\binom{r_-}{2}
                                      +\binom{r_-}{3}\right\},
\label{ms:discrete-curvature}
\end{align}
where $r_+=\max(r,0)$ and $r_-=\max(-r,0)$. If $|u|\le1/2$ and
$t=\E V_k^2$, then
\[
 \E R_+\le\frac{t}{(\zeta-u)^2},\qquad
 \E\binom{R_-}{2}\le\frac{t}{[\zeta+\min(u,\zeta/3)]^2}.
\]
If additionally $|U_0|,|U_1|\le A_U$, then
\[
 \E\binom{R_-}{3}\le
           \frac{t\sqrt{A_U^2+3t}}{(2\zeta-1/2)^3}.
\]
\end{lemma}
\begin{proof}
For $0\le j\le n$, the binomial difference identity is
\[
 \pi_j-\pi_{j-1}
    =\pi_j\frac{(n+1)p-j}{p(n-j+1)}.
\]
When its numerator is positive, $p(n-j+1)\ge(n+1)v$.
Thus its positive part is bounded by
$m_B((n+1)p-j)_+/((n+1)v)$; zero extension gives the same upper
bound outside the legal range. For $r>0$, double summation gives weights
$r-i+1$ on these differences at $j=k+i$. Since $\chi<2$, only $i=1$
can contribute positively, giving $(\chi-1)_+r$.
For $r=-h<0$, the corresponding weights are $h-1-i$ on the differences
at $j=k-i$, $0\le i\le h-2$. Use $(\chi+i)_+\le\chi_++i$ and
sum the weights to obtain $\chi_+\binom h2+\binom h3$.

On $R=r\ge1$, crossing monotonicity gives $-V_k\ge \zeta r-u$.
Since $(\zeta r-u)^2\ge r(\zeta-u)^2$, this proves the first moment bound.
On $R=-h\le-2$, the count immediately above $K$ is outside the closed
event, giving the strict inequality $V_k>\zeta(h-1)+u$. Put $j=h-1$.
For $u\le \zeta/3$, differentiation shows that
$j(j+1)/(2(\zeta j+u)^2)$ decreases for $j\ge1$: the derivative has the sign
of $j(2u-\zeta)+u\le0$. Its maximum is $1/(\zeta+u)^2$ at $j=1$.
For $u>\zeta/3$, replace the crossing threshold by the smaller one with
$u=\zeta/3$. This proves the second bound.

For the third, use $u\ge-1/2$ and $j\ge2$. The derivative of
$(j^3-j)/(6(\zeta j-1/2)^3)$ has the sign of
$2\zeta j-(3j^2-1)/2<0$, because $\zeta\le1$. Its maximum at $j=2$ is
$(2\zeta-1/2)^{-3}$. Hence
$\binom{R_-}{3}\le(V_k)_+^3/(2\zeta-1/2)^3$.
Independence and centering give
$\E V_k^4\le A_U^2t+3t^2$. Finally,
$\E(V_k)_+^3\le\sqrt{\E V_k^2\,\E V_k^4}
\le t\sqrt{A_U^2+3t}$ by Cauchy--Schwarz.
\end{proof}

Here $\chi=p-H/\ell+u$. Thus an upper bound
$b=p_+-\min(H/\ell)$ yields precisely the phase function
\[
 \frac{(b+u-1)_+}{(\zeta-u)^2}
 +\frac{(b+u)_+}{[\zeta+\min(u,\zeta/3)]^2}.
\]
We maximize this function in Section~\ref{ms:curvature}.
The conditions $\chi<2$ and $1-c_v\Delta/N>0$, and any normalization
such as $\sqrt n\,m_B\le1$, must be checked separately.
The marked-crossing estimates needed to bound the remaining linear term
$\pi_k\E R$ in \eqref{ms:whole-count-cutoff} are given in Section~\ref{ms:fourier:local}; they form a separate part
of the local comparison.
